\section{Availability, reproducibility, and intended use}

\texttt{voiage} is intended to support research in health technology
assessment, clinical-trial design, public policy, environmental
decision-making, and related areas of uncertainty quantification. The
repository contains deterministic analytical and integration fixtures,
tutorials, cross-language contracts, and representative workflows for these
use cases.

Version 1.0.0 is available under the Apache License 2.0 from
\url{https://github.com/edithatogo/voiage/releases/tag/v1.0.0} and
\url{https://pypi.org/project/voiage/1.0.0/}.
The repository records citation metadata, supported Python versions, the Rust
minimum supported version, locked dependencies, release checksums, and
machine-readable provenance. The documentation site is built from the same
repository using Astro/Starlight. Examples and normative fixtures use synthetic
or explicitly governed inputs so that the package can be evaluated without
access to confidential decision models.

The project applies the principle that research software should be findable,
accessible, interoperable, and reusable where technically
possible~\cite{wilkinson2016fair}. In practice, this means stable identifiers for
schemas, portable fixtures, explicit licenses, documented installation paths,
and result envelopes that can be inspected outside the originating language.
It does not mean that every experimental method or external registry has the
same maturity.

Potential applications include comparing research designs, prioritizing
parameters for further study, assessing decision uncertainty across
subgroups, and embedding VOI results in larger evidence pipelines. These are
intended uses, not evidence of adoption or policy impact. Publicly attributable
case studies can be added in a later manuscript revision when such evidence is
available.
